\chapter{研究现状}\label{chap:related_work}

% 本章 TikZ 图形统一样式定义
\tikzset{
    % 基础模块样式
    module/.style={
        draw, rounded corners=2pt, minimum height=0.7cm,
        font=\small, align=center, thick,
        fill=#1!12, draw=#1!70!black
    },
    module/.default=blue,
    % 箭头样式
    arrow/.style={-{Stealth[length=5pt]}, thick},
    arrow2/.style={-{Stealth[length=4pt]}, semithick},
    dasharrow/.style={-{Stealth[length=5pt]}, thick, dashed},
    % 大括号标注
    brace/.style={decorate, decoration={brace, amplitude=5pt, mirror, raise=2pt}},
    braceup/.style={decorate, decoration={brace, amplitude=5pt, raise=2pt}},
    % 标签
    lbl/.style={font=\footnotesize, align=center},
    slbl/.style={font=\scriptsize, align=center},
    % 虚线框
    dbox/.style={draw, dashed, rounded corners=3pt, inner sep=6pt, #1},
    dbox/.default={gray!60},
}

语音到语音生成是本文各章工作的共同研究主题。从基于编码器-解码器架构的语音翻译，到基于大语言模型的语音对话，本文的五项工作分别从不同角度推动了语音到语音生成技术的发展。为全面理解这一领域的研究脉络，本章将系统性地梳理语音到语音生成的相关研究现状。


%===========================================================================
% 2.1 语音到语音生成的建模范式
%===========================================================================
\section{语音到语音生成的建模范式}\label{sec:rw_paradigm}

语音到语音生成（Speech-to-Speech Generation, S2S）是一类以语音为输入、以语音为输出的序列到序列生成任务。典型的语音到语音生成任务包括语音到语音翻译（将源语言语音转换为目标语言语音）和语音对话（根据用户的语音指令生成语音回复）等。从系统层面来看，如图\ref{fig:rw_paradigm}所示，语音到语音生成存在三种主要的建模范式，分别对应不同的模块化程度和端到端程度。

% 图：建模范式对比图
\begin{figure}[htbp]
\centering
\begin{tikzpicture}[node distance=0.5cm and 0.5cm]
    \def\xleft{-0.4}
    \def\xright{11.4}
    \def\xcenter{5.5}

    % === (a) 三阶段级联系统 ===
    \def\ya{0}
    \node[module=green, minimum width=2.0cm, minimum height=0.9cm] (asr) at (2.0, \ya) {语音识别};
    \node[module=orange, minimum width=2.4cm, minimum height=0.9cm] (mt) at (\xcenter, \ya) {机器翻译/\\[-2pt]大语言模型};
    \node[module=red, minimum width=2.0cm, minimum height=0.9cm] (tts) at (9.0, \ya) {语音合成};
    \node[font=\small] (sin1) at (\xleft, \ya) {输入语音};
    \node[font=\small] (sout1) at (\xright, \ya) {输出语音};
    \draw[arrow] (sin1) -- (asr);
    \draw[arrow] (asr) -- node[above, yshift=1pt, font=\scriptsize]{转写文本} (mt);
    \draw[arrow] (mt) -- node[above, yshift=1pt, font=\scriptsize]{目标文本} (tts);
    \draw[arrow] (tts) -- (sout1);
    \node[font=\small\bfseries] at (\xcenter, \ya-1.0) {(a) 三阶段级联系统};

    % === (b) 两阶段级联系统 ===
    \def\yb{-2.5}
    \node[module=teal, minimum width=5.7cm, minimum height=0.9cm] (s2t) at (3.85, \yb) {语音到文本模型};
    \node[module=red, minimum width=2.0cm, minimum height=0.9cm] (tts2) at (9.0, \yb) {语音合成};
    \node[font=\small] (sin2) at (\xleft, \yb) {输入语音};
    \node[font=\small] (sout2) at (\xright, \yb) {输出语音};
    \draw[arrow] (sin2) -- (s2t);
    \draw[arrow] (s2t) -- node[above, yshift=1pt, font=\scriptsize]{目标文本} (tts2);
    \draw[arrow] (tts2) -- (sout2);
    \node[font=\small\bfseries] at (\xcenter, \yb-0.8) {(b) 两阶段级联系统};

    % === (c) 端到端模型 ===
    \def\yc{-4.5}
    \node[module=violet, minimum width=9.0cm, minimum height=0.9cm] (e2e) at (\xcenter, \yc) {端到端语音到语音模型};
    \node[font=\small] (sin3) at (\xleft, \yc) {输入语音};
    \node[font=\small] (sout3) at (\xright, \yc) {输出语音};
    \draw[arrow] (sin3) -- (e2e);
    \draw[arrow] (e2e) -- (sout3);
    \node[font=\small\bfseries] at (\xcenter, \yc-0.8) {(c) 端到端模型};
\end{tikzpicture}
\bicaption{\enspace 语音到语音生成的三种建模范式}{\enspace Three modeling paradigms for speech-to-speech generation}
\label{fig:rw_paradigm}
\end{figure}

第一种是三阶段级联范式。该范式将语音到语音生成过程拆分为三个独立模块的串联：自动语音识别（ASR）模块将输入语音转录为源语言文本，机器翻译（MT）或大语言模型（LLM）模块对文本进行处理或翻译，语音合成（TTS）模块将目标文本转换为输出语音。早期的语音翻译系统普遍采用这一范式\citep{lavie1997janus,ney1999speech,nakamura2006atr}。级联范式的优势在于模块化程度高，各模块可以独立开发和优化，并能复用成熟的预训练模型。然而，该范式存在固有的缺陷：上游模块的识别错误会逐级累积并放大，导致最终输出质量下降；多个模块的串行推理带来较高的延迟；语音中的副语言信息（如语调、情感、韵律）在文本中间表示中丢失，难以在目标语音中保持。

第二种是两阶段级联范式。该范式将语音识别与文本处理整合为一个端到端的语音到文本（Speech-to-Text, S2T）模型，再将其输出送入语音合成模块生成目标语音。相比三阶段级联范式，两阶段级联范式将语音识别和机器翻译（或文本处理）合并为单一的语音到文本模型进行联合优化，减少了模块间的信息损失和错误累积。同时，由于S2T模型可以直接学习从源语音到目标文本的映射，避免了中间转写步骤引入的噪声。两阶段级联范式本质上仍是级联架构——S2T模型和TTS模型作为两个独立模块串联执行，各模块可以独立训练和替换。这种模块化的特性使得两阶段级联系统能够直接复用S2T和TTS领域已有的高质量预训练模型，是当前实际部署中广泛采用的方案。

第三种是端到端范式。该范式通过一个统一模型直接完成从输入语音到输出语音的映射，无需显式的文本中间表示。端到端范式在理论上能够消除级联系统的误差累积和信息丢失问题，并降低多模块串行推理带来的延迟。更重要的是，由于不依赖文本作为中间表示，端到端模型有可能保留语音中的副语言信息（如说话人音色、语调、情感和韵律等），实现更加自然的语音到语音转换。然而，端到端范式也面临着训练数据需求大、建模难度高等挑战，尤其是直接从语音到语音的映射需要模型同时学习语音理解和语音生成两项能力。尽管如此，端到端范式仍是当前语音到语音生成研究追求的核心目标。

在模型架构层面，上述建模范式的演进与序列到序列生成领域的架构发展密切相关。2017年，\citet{transformer}提出的Transformer模型以编码器-解码器（Encoder-Decoder）架构为基础，依靠自注意力机制取代了传统的循环结构，在机器翻译等序列到序列任务上取得了突破性的成果。此后的数年间，编码器-解码器架构成为序列到序列生成任务的主流范式，广泛应用于语音识别、语音翻译和语音到语音翻译等任务。2018年起，GPT\citep{gpt1}等工作证明了仅解码器（Decoder-only）架构在语言建模任务上的强大能力。随着模型规模和训练数据的持续增长，以LLaMA\citep{touvron2023llama}、Qwen\citep{yang2024qwen2}为代表的大语言模型广泛采用仅解码器架构，并展现出卓越的语言理解和生成能力。在此背景下，语音到语音生成领域也逐步从编码器-解码器架构向基于大语言模型的仅解码器架构过渡。

如图\ref{fig:rw_overview}所示，本章将沿着从语音翻译到语音对话的演进主线展开综述。语音翻译方向的研究主要基于编码器-解码器架构，涵盖语音到文本翻译和语音到语音翻译；语音对话方向的研究主要基于大语言模型，涵盖语音理解和语音对话。在此之前，第\ref{sec:rw_repr}节先介绍语音表征的相关研究，为后续内容提供前置知识。最后，第\ref{sec:rw_analysis}节对研究现状进行综合分析，指出当前研究中存在的关键问题，引出本文的研究内容。

% 图：本章研究现状分类概览（树状图）
\begin{figure}[htbp]
\centering
\begin{tikzpicture}[
    node distance=0.6cm and 1.2cm,
    box/.style={draw, rounded corners=3pt, minimum height=0.8cm,
        font=\small, align=center, thick, fill=#1!12, draw=#1!70!black},
    box/.default=blue,
]
    \node[box=violet, minimum width=3.0cm] (root) at (5.5, 0) {语音到语音生成};
    \node[box=blue, minimum width=3.2cm] (encdec) at (2.5, -1.5) {语音翻译\\[-2pt]{\scriptsize 编码器-解码器（\S\ref{sec:rw_encdec_all}）}};
    \node[box=orange, minimum width=3.2cm] (deconly) at (8.5, -1.5) {语音对话\\[-2pt]{\scriptsize 大语言模型（\S\ref{sec:rw_deconly_all}）}};
    \node[box=blue, inner xsep=6pt] (s2tt) at (0.0, -3.8) {语音到文本\\[-2pt]{\scriptsize（\S\ref{sec:rw_encdec_s2tt} 语音到文本翻译）}};
    \node[box=blue, inner xsep=6pt] (s2st) at (3.7, -3.8) {语音到语音\\[-2pt]{\scriptsize（\S\ref{sec:rw_encdec_s2st} 语音到语音翻译）}};
    \node[box=orange, inner xsep=6pt] (s2t_llm) at (7.3, -3.8) {语音到文本\\[-2pt]{\scriptsize（\S\ref{sec:rw_deconly_s2t} 语音理解大模型）}};
    \node[box=orange, inner xsep=6pt] (s2s_llm) at (11.0, -3.8) {语音到语音\\[-2pt]{\scriptsize（\S\ref{sec:rw_deconly_s2s} 语音对话大模型）}};
    \draw[arrow] (root) -- (encdec);
    \draw[arrow] (root) -- (deconly);
    \draw[arrow] (encdec) -- (s2tt);
    \draw[arrow] (encdec) -- (s2st);
    \draw[arrow] (deconly) -- (s2t_llm);
    \draw[arrow] (deconly) -- (s2s_llm);
\end{tikzpicture}
\bicaption{\enspace 语音到语音生成研究现状的分类概览}{\enspace A taxonomy of speech-to-speech generation research}
\label{fig:rw_overview}
\end{figure}


%===========================================================================
% 2.2 语音表征
%===========================================================================
\section{语音表征}\label{sec:rw_repr}

如第\ref{chap:introduction}章所述，语音到语音生成的完整流程为：将输入语音波形转换为输入语音表征，由核心模型将输入语音表征映射为输出语音表征，再通过声码器将输出语音表征还原为语音波形。原始语音波形是以数万赫兹采样率录制的时域信号，信息密度低且冗余度高，难以直接用于序列到序列的建模。因此，将语音波形转换为更紧凑、更结构化的中间表征，是构建语音到语音生成系统的基础。本节将分别介绍用于理解的语音表征和用于生成的语音表征的相关研究。

\subsection{用于理解的语音表征}\label{sec:rw_repr_understand}

用于理解的语音表征旨在将原始语音波形转化为可被下游模型处理的特征。根据表征的层次不同，可分为底层声学特征和高层语音表征两类。底层声学特征是从波形直接提取的信号级特征（如滤波器组特征），不涉及模型学习，可直接作为模型的输入。在此之上，语音编码器通过对底层特征进行上下文建模，提取出包含更丰富信息的高层语音表征。当语音编码器的参数固定时，它可以被视为一个离线的特征提取器，其输出同样可以作为下游模型的输入表征；当编码器作为模型的一部分参与端到端训练时，高层表征则随训练动态更新。

\textbf{底层声学特征}\quad 最常用的底层声学特征是滤波器组（Filterbank, Fbank）特征，它将语音波形经短时傅里叶变换和梅尔滤波器组处理，得到一组反映频谱特性的连续特征。梅尔谱（Mel Spectrogram）是Fbank特征的一种常见形式，以时间-频率的二维矩阵表示语音信号的频谱包络信息，梅尔标度与人耳的感知特性近似匹配。在预训练语音编码器出现之前，Fbank特征是语音识别、语音翻译等任务的标准输入，模型通过语音编码器与下游任务端到端联合训练来学习从底层特征到高层表征的映射。

\textbf{高层语音表征}\quad 语音编码器通过对底层声学特征进行深层编码，提取出包含上下文信息的高层语音表征。在编码器架构方面，Conformer\citep{conformer}通过在Transformer自注意力模块中引入卷积模块，同时捕获语音信号的局部特征和全局依赖，成为语音识别和语音翻译任务中广泛使用的编码器架构。E-Branchformer\citep{kim2023ebranchformer}提出了增强合并的分支结构，Zipformer\citep{yao2024zipformer}提出了多速率U-Net结构，FastConformer\citep{rekesh2023fastconformer}将全局注意力替换为线性可扩展的局部注意力，它们均在Conformer的基础上进一步提升了编码效率或精度。此外，Paraformer\citep{gao2022paraformer}探索了非自回归的端到端语音识别架构，通过连续积分与激发（Continuous Integrate-and-Fire, CIF）\citep{dong2020cif}机制实现并行解码。

在预训练方法方面，自监督和监督两种范式均取得了重要进展。在自监督预训练中，CPC\citep{oord2018cpc}和wav2vec\citep{schneider2019wav2vec}最早通过对比预测学习语音的上下文表征。Wav2Vec 2.0\citep{baevski2020wav2vec}统一了对比学习与掩码预测的框架，为后续工作奠定了基础。HuBERT\citep{hubert}采用离线聚类生成伪标签并通过掩码预测进行训练，学习到的表征在多项下游任务上表现优异，且其中间层特征经聚类后可作为离散语音单元广泛应用于语音到语音翻译。WavLM\citep{chen2022wavlm}在掩码预测的基础上引入去噪任务，W2V-BERT\citep{chung2021w2vbert}将对比学习与掩码语言建模相结合。在模型规模扩展方面，MMS\citep{pratap2023mms}将自监督预训练扩展到超过1100种语言的50万小时数据，Google USM\citep{zhang2023usm}进一步扩展到1200万小时、300余种语言。在音频理解方面，SSAST\citep{gong2022ssast}、Audio-MAE\citep{huang2022audiomae}和BEATs\citep{chen2023beats}将视觉领域的自监督方法迁移到音频频谱图上，取得了优异的性能。在监督预训练中，Whisper\citep{whisper}在68万小时的弱监督语音识别数据上进行训练，覆盖98种语言，其编码器提取的语音表征具有良好的跨语言和跨任务迁移能力，被广泛用作后续语音大模型的前端编码器。OWSM\citep{peng2023owsm}以开源方式复现了Whisper的多任务多语言训练范式，SenseVoice\citep{an2024funaudiollm}则面向50余种语言构建了多任务语音理解模型。


\subsection{用于生成的语音表征}\label{sec:rw_repr_gen}

语音生成模型需要一种目标表征来定义生成的输出。根据表征形式的不同，用于生成的语音表征可分为连续表征（梅尔谱和VAE潜在特征）和离散表征（单码本和多码本离散词元）两大类。

\textbf{连续频谱特征}\quad 梅尔谱（Mel Spectrogram）是语音生成中最基础的表征形式。它将原始波形通过短时傅里叶变换转为频谱，再经梅尔滤波器组映射到梅尔标度上，得到一个时间-频率的二维矩阵。梅尔谱保留了语音信号的频谱包络信息，能够反映音素、音高和音色等声学特性，同时梅尔标度与人耳的感知特性近似匹配，使得基于梅尔谱的生成结果具有较好的感知质量。梅尔谱是早期端到端语音合成模型（如Tacotron\citep{wang2017tacotron}、Tacotron 2\citep{shen2018tacotron2}和FastSpeech\citep{ren2019fastspeech}）和语音到语音翻译模型（如Translatotron\citep{translatotron}）的默认生成目标。由于梅尔谱本身不包含相位信息，无法直接还原为语音波形，因此需要通过声码器进行波形重建。Griffin-Lim\citep{griffin1984signal}算法是最早采用的相位估计方法，但生成质量有限。随后，基于神经网络的声码器逐渐成为主流，WaveNet\citep{oord2016wavenet}以自回归方式逐采样点生成波形，生成质量极高但推理速度慢。HiFi-GAN\citep{hifigan}采用生成对抗网络架构，通过多周期判别器和多尺度判别器实现了高保真度和快速推理的平衡，成为目前最广泛使用的声码器之一。

\textbf{连续VAE潜在特征}\quad 近年来，基于变分自编码器（Variational Autoencoder, VAE）的连续潜在表征逐渐成为语音和音频生成的新范式。VAE编码器将波形或频谱压缩到低维潜在空间中，解码器则将潜在表示还原为波形。相比梅尔谱，VAE潜在特征维度更低、更紧凑，降低了下游生成模型的建模负担，且其连续且平滑的潜在空间天然适合与扩散模型结合进行高质量生成。AudioLDM\citep{liu2023audioldm}借鉴图像生成中潜在扩散模型（Latent Diffusion Model）的思想，在梅尔谱的VAE潜在空间上运行扩散模型，以文本提示引导音频生成，在文本到音频生成任务上取得了优异的效果。NaturalSpeech 2\citep{shen2024naturalspeech2}将扩散模型应用于神经编解码器的连续潜在向量，通过时长和音高预测器实现了零样本语音合成。NaturalSpeech 3\citep{ju2024naturalspeech3}进一步提出分解式扩散方法，将语音分解为内容、韵律、音色和声学细节等子空间，分别建模以提升合成的可控性和质量。Stable Audio\citep{evans2024stableaudio}采用全卷积VAE与扩散Transformer的组合，以时间和文本条件引导长时音频的生成，支持高达44.1kHz采样率的高保真音频生成。Ming-UniAudio\citep{liang2025mingtok}将VAE连续潜在表征作为语音大语言模型的统一表示，使单一模型能够同时完成语音理解、生成和编辑任务。

\textbf{离散单码本词元}\quad 单码本表示将每帧语音映射为一个离散词元，一段语音在给定帧率下产生一维的词元序列。单码本离散词元主要有两种获取方式。第一种是自监督模型加聚类的方式：利用HuBERT等自监督模型提取中间层特征，再经K-Means聚类产生语义离散单元，该方法以50 Hz帧率广泛应用于语音到语音翻译中的离散单元预测。第二种是基于单码本神经编解码器的方式：CosyVoice\citep{du2024cosyvoice}的S3 Tokenizer在预训练语音识别模型的编码器中插入向量量化层，通过语音识别损失进行监督训练，产生语义丰富的单码本词元。CosyVoice 2\citep{du2024cosyvoice2}进一步采用有限标量量化（Finite Scalar Quantization, FSQ）替代向量量化，实现了接近100\%的码本利用率。GLM-4-Voice\citep{glm4voice}在Whisper编码器中插入向量量化瓶颈，以12.5 Hz的超低帧率实现了语义词元的高效提取。WavTokenizer\citep{ji2024wavtokenizer}和BigCodec\citep{xin2024bigcodec}则从编解码器架构出发设计单码本方案，分别在重建质量和极低比特率上取得了优异的表现。Single-Codec\citep{li2024singlecodec}通过解耦说话人嵌入和语音内容实现了单码本下的高质量语音重建。RepCodec\citep{huang2023repcodec}则从语音表征压缩的角度出发，通过对预训练语音模型的表征进行向量量化，构建了面向语音分词的表征编解码器。

\textbf{离散多码本词元}\quad 多码本表示通过残差向量量化（Residual Vector Quantization, RVQ）将每帧语音映射为多个离散词元，通常第一层编码语义信息，后续层逐步编码声学细节。SoundStream\citep{zeghidour2021soundstream}提出了基于RVQ和对抗训练的神经音频编解码器框架。EnCodec\citep{defossez2022encodec}在此基础上实现了多码率的高保真音频压缩，被VALL-E\citep{wang-etal-2023-valle}和AudioLM\citep{borsos2023audiolm}等语音语言模型广泛采用。DAC\citep{kumar2023dac}通过改进的RVQGAN架构实现了44.1kHz音频的约90倍压缩。SpeechTokenizer\citep{zhang2024speechtokenizer}通过HuBERT蒸馏将RVQ第一层解耦为语义层，后续层编码声学信息，实现了语义和声学信息的显式分离。SNAC\citep{siuzdak2024snac}提出了多尺度RVQ架构，以12/24/48 Hz三个层级的不同帧率编码语音信息。Mimi\citep{kyutai2024moshi}同样实现了语义与声学的解耦，以12.5 Hz的低帧率支持实时对话建模。X-Codec\citep{ye2024xcodec}则通过语义增强策略提升了编解码器对语义信息的保留能力。SemantiCodec\citep{liu2024semanticodec}将语义信息编码与声学信息编码分离，在极低比特率（约0.31 kbps）下实现了通用音频的高质量编解码。XY-Tokenizer\citep{xiao2025xytokenizer}针对低比特率编解码器中语义信息和声学信息之间的冲突问题，提出了语义-声学解耦的双码本架构，在低比特率条件下同时保持了语义完整性和声学质量。


%===========================================================================
% 2.3 基于编码器-解码器的语音到语音生成
%===========================================================================
\section{基于编码器-解码器的语音翻译}\label{sec:rw_encdec_all}

编码器-解码器架构是端到端语音翻译研究的核心架构。本节首先介绍编码器-解码器架构的基础知识，随后分别介绍基于该架构的语音到文本翻译和语音到语音翻译的相关研究。

\subsection{编码器-解码器架构基础}\label{sec:rw_encdec}

编码器-解码器架构是基于Transformer的序列到序列建模的经典范式，也是早期语音到语音生成研究的核心架构。\citet{transformer}于2017年提出的Transformer模型彻底摒弃了传统的循环和卷积结构，完全依赖自注意力机制来捕捉序列内部任意位置之间的依赖关系，并凭借高度并行化的结构大幅提升了训练效率。

如图\ref{fig:rw_transformer}所示，Transformer编码器-解码器架构由编码器和解码器两部分组成，二者均由若干层相同结构的层堆叠而成。编码器负责对源序列$X=(x_1,\ldots,x_N)$进行编码，得到连续向量表示$Z=(z_1,\ldots,z_N)$。每个编码器层包含两个子模块：双向自注意力模块和前馈神经网络。双向自注意力允许序列中的每个位置关注所有其他位置，从而捕捉完整的上下文信息。

% 图：Transformer编码器-解码器架构示意图
\begin{figure}[htbp]
\centering
\begin{tikzpicture}[node distance=0.35cm,
    block/.style={draw, rounded corners=2pt, minimum width=2.4cm, minimum height=0.6cm,
        font=\small, align=center, thick, fill=#1!12, draw=#1!70!black},
    block/.default=blue]

    % 编码器（自底向上）
    \node[block=green, minimum width=2.8cm] (enc_input) at (-2.5, 0) {输入嵌入 + 位置编码};
    \node[block=blue, above=of enc_input, minimum width=2.8cm] (enc_sa) {双向自注意力};
    \node[block=blue, above=of enc_sa, minimum width=2.8cm] (enc_ffn) {前馈网络};
    \draw[arrow2] (enc_input) -- (enc_sa);
    \draw[arrow2] (enc_sa) -- (enc_ffn);
    \begin{scope}[on background layer]
        \node[dbox=blue!40, fit=(enc_sa)(enc_ffn), inner sep=5pt] (enc_layer) {};
    \end{scope}
    \node[right=-2pt of enc_layer.east, anchor=west, font=\scriptsize, blue!70] {$\times N$};
    % 编码器输出
    \node[above=0.35cm of enc_ffn, font=\small, fill=white, inner sep=2pt] (enc_out) {编码器输出 $Z$};
    \draw[arrow2] (enc_ffn) -- (enc_out);

    % 解码器（自底向上）
    \node[block=orange, minimum width=2.8cm] (dec_input) at (3.0, 0) {输出嵌入 + 位置编码};
    \node[block=red, above=of dec_input, minimum width=2.8cm] (dec_sa) {因果自注意力};
    \node[block=violet, above=of dec_sa, minimum width=2.8cm] (dec_ca) {交叉注意力};
    \node[block=red, above=of dec_ca, minimum width=2.8cm] (dec_ffn) {前馈网络};
    \node[block=gray, above=of dec_ffn, minimum width=2.8cm] (dec_linear) {线性层 + Softmax};
    \draw[arrow2] (dec_input) -- (dec_sa);
    \draw[arrow2] (dec_sa) -- (dec_ca);
    \draw[arrow2] (dec_ca) -- (dec_ffn);
    \draw[arrow2] (dec_ffn) -- (dec_linear);
    \begin{scope}[on background layer]
        \node[dbox=red!40, fit=(dec_sa)(dec_ca)(dec_ffn), inner sep=5pt] (dec_layer) {};
    \end{scope}
    \node[right=-2pt of dec_layer.east, anchor=west, font=\scriptsize, red!70] {$\times N$};
    % 输出
    \node[above=0.2cm of dec_linear, font=\small] (output) {输出};
    \draw[arrow2] (dec_linear) -- (output);

    % 交叉注意力连接
    \draw[arrow2, violet!70] (enc_out.east) -| ([xshift=-0.4cm]dec_ca.west) -- (dec_ca.west);

    % 编码器/解码器大虚线框 + 底部标题
    \begin{scope}[on background layer]
        \node[draw, dashed, rounded corners=4pt, gray!50,
            fit=(enc_input)(enc_out)(enc_layer),
            inner xsep=8pt, inner ysep=10pt] (enc_box) {};
        \node[draw, dashed, rounded corners=4pt, gray!50,
            fit=(dec_input)(output)(dec_layer),
            inner xsep=8pt, inner ysep=10pt] (dec_box) {};
    \end{scope}
    \node[font=\small\bfseries, gray!80, below=3pt of enc_box] {编码器};
    \node[font=\small\bfseries, gray!80, below=3pt of dec_box] {解码器};
\end{tikzpicture}
\bicaption{\enspace Transformer编码器-解码器架构}{\enspace Transformer encoder-decoder architecture}
\label{fig:rw_transformer}
\end{figure}

解码器负责根据编码器的输出表示$Z$，以自回归方式逐步生成目标序列$Y=(y_1,\ldots,y_M)$。每个解码器层包含三个子模块：因果自注意力模块、交叉注意力模块和前馈神经网络。因果自注意力通过掩码确保在生成位置$i$的预测时仅能关注位置$i$及其之前的信息，维持了生成过程的因果性。交叉注意力是连接编码器和解码器的桥梁，其查询来自解码器，键和值来自编码器输出，使解码器在生成每个目标词时能够选择性地关注源序列的不同部分。

Transformer的核心组件包括多头注意力机制、前馈神经网络和位置编码。多头注意力基于缩放点积注意力，将查询、键、值通过多组独立的线性变换投影后分别计算注意力，再将各组输出拼接并线性变换，使模型能够在不同的表示子空间中关注不同位置的信息。前馈神经网络为每个位置的表示提供额外的非线性变换。由于自注意力机制本身不包含位置信息，Transformer通过位置编码为输入序列注入位置信息。整个模型通过最小化目标序列的负对数似然进行训练：
\begin{equation}
\mathcal{L}(\theta)=-\sum_{i=1}^{M}\log P_\theta(y_i|X,Y_{<i}),
\end{equation}
其中$\theta$为模型参数，$Y_{<i}$为位置$i$之前的目标序列前缀。

编码器-解码器架构适用于需要将一个序列映射到另一个序列的场景，在机器翻译、语音识别和语音翻译等任务中得到了广泛应用。在语音到语音生成领域，该架构构成了大量端到端语音翻译和语音到语音翻译工作的基础。


\subsection{语音到文本翻译}\label{sec:rw_encdec_s2tt}

基于编码器-解码器的语音到文本生成是语音到语音生成的重要前置任务，也是本文第三章和第四章工作的基础。端到端语音到文本翻译（Speech-to-Text Translation, S2TT）直接将源语言语音信号转换为目标语言文本，相比传统的ASR加MT级联系统，在减少错误累积和降低延迟方面具有显著优势。然而，S2TT面临的核心挑战在于高质量语音-文本平行数据的匮乏，以及语音与文本之间的模态鸿沟（Modality Gap）。围绕这些挑战，研究者从模型结构和训练方法两个维度展开了大量研究。按照模型结构，如图\ref{fig:rw_s2tt}所示，现有S2TT方法可分为三类：基础编码器-解码器、并联双编码器和串联双编码器。

% 图：语音到文本翻译模型结构对比图
% 设计原则：三子图整体垂直居中；箭头视觉长度统一
% 关键：TikZ从节点边框画箭头，文本标签比方框矮，需加大标签↔方框的中心距来补偿
\begin{figure}[htbp]
\centering
\begin{tikzpicture}[node distance=0.3cm and 0.3cm,
    enc/.style={draw, rounded corners=2pt, minimum height=0.7cm, minimum width=2.4cm,
        font=\small, align=center, thick, fill=#1!12, draw=#1!70!black},
    enc/.default=blue,
]
    % 间距参数（控制箭头视觉长度）
    % 方框高0.7cm，标签高约0.25cm；TikZ从边框画箭头
    % 方框↔方框：视觉箭头 = 中心距 - 0.7
    % 标签↔方框：视觉箭头 = 中心距 - 0.35 - 0.125 ≈ 中心距 - 0.475
    % 要使视觉相等，标签的中心距需比方框的中心距多约0.225
    \def\boxgap{0.4}     % 方框之间的视觉箭头长度
    \def\lblgap{0.63}    % 标签↔方框的中心距偏移（补偿标签较矮）
    % 验证：方框视觉 = 0.4；标签视觉 ≈ 0.63 - 0.475 ≈ 0.4 — 匹配 ✓
    \def\titleGap{0.7}   % 底部标签与标题的间距

    % === (c) 串联双编码器 — 最高的子图，先定位以确定全局范围 ===
    \def\xc{10.0}
    \pgfmathsetmacro{\cMid}{0}
    \pgfmathsetmacro{\cTenc}{\cMid}                        % 文本编码器（中心）
    \pgfmathsetmacro{\cSenc}{\cMid-0.7-\boxgap}            % 语音编码器
    \pgfmathsetmacro{\cDec}{\cMid+0.7+\boxgap}             % 文本解码器
    \pgfmathsetmacro{\cIn}{\cSenc-0.35-\lblgap}             % 输入标签
    \pgfmathsetmacro{\cOut}{\cDec+0.35+\lblgap}             % 输出标签
    \pgfmathsetmacro{\titleY}{\cIn-\titleGap}               % 标题（三图共用）

    \node[enc=green] (c_senc) at (\xc, \cSenc) {语音编码器};
    \node[enc=orange] (c_tenc) at (\xc, \cTenc) {文本编码器};
    \node[enc=red] (c_dec) at (\xc, \cDec) {文本解码器};
    \node[font=\scriptsize] (c_in) at (\xc, \cIn) {源语言语音};
    \node[font=\scriptsize] (c_out) at (\xc, \cOut) {目标语言文本};
    \draw[arrow2] (c_in) -- (c_senc);
    \draw[arrow2] (c_senc) -- (c_tenc);
    \draw[arrow2] (c_tenc) -- (c_dec);
    \draw[arrow2] (c_dec) -- (c_out);
    \node[font=\small\bfseries] at (\xc, \titleY) {(c) 串联双编码器};

    % === (a) 单编码器 — 垂直居中 ===
    \def\xa{0}
    \pgfmathsetmacro{\aMid}{\cMid}
    \pgfmathsetmacro{\aEnc}{\aMid-0.35-\boxgap/2}          % 语音编码器
    \pgfmathsetmacro{\aDec}{\aMid+0.35+\boxgap/2}          % 文本解码器
    \pgfmathsetmacro{\aIn}{\aEnc-0.35-\lblgap}              % 输入标签
    \pgfmathsetmacro{\aOut}{\aDec+0.35+\lblgap}             % 输出标签

    \node[enc=green] (a_senc) at (\xa, \aEnc) {语音编码器};
    \node[enc=red] (a_dec) at (\xa, \aDec) {文本解码器};
    \node[font=\scriptsize] (a_in) at (\xa, \aIn) {源语言语音};
    \node[font=\scriptsize] (a_out) at (\xa, \aOut) {目标语言文本};
    \draw[arrow2] (a_in) -- (a_senc);
    \draw[arrow2] (a_senc) -- (a_dec);
    \draw[arrow2] (a_dec) -- (a_out);
    \node[font=\small\bfseries] at (\xa, \titleY) {(a) 单编码器};

    % === (b) 并联双编码器 — 与(a)相同高度，垂直居中 ===
    \def\xb{5.0}
    \node[enc=green] (b_senc) at (\xb-1.4, \aEnc) {语音编码器};
    \node[enc=orange] (b_tenc) at (\xb+1.4, \aEnc) {文本编码器};
    \node[enc=red] (b_dec) at (\xb, \aDec) {文本解码器};
    \node[font=\scriptsize] (b_in1) at (\xb-1.4, \aIn) {源语言语音};
    \node[font=\scriptsize] (b_in2) at (\xb+1.4, \aIn) {源语言文本};
    \node[font=\scriptsize] (b_out) at (\xb, \aOut) {目标语言文本};
    \draw[arrow2] (b_in1) -- (b_senc);
    \draw[arrow2] (b_in2) -- (b_tenc);
    \draw[arrow2] (b_senc.north) -- ([xshift=-0.6cm]b_dec.south);
    \draw[arrow2] (b_tenc.north) -- ([xshift=0.6cm]b_dec.south);
    \draw[arrow2] (b_dec) -- (b_out);
    \node[font=\small\bfseries] at (\xb, \titleY) {(b) 并联双编码器};
\end{tikzpicture}
\bicaption{\enspace 语音到文本翻译的三种模型结构}{\enspace Three model architectures for speech-to-text translation}
\label{fig:rw_s2tt}
\end{figure}

\subsubsection{基础编码器-解码器}\label{sec:rw_s2tt_basic}

基础编码器-解码器是最直接的端到端S2TT结构，由语音编码器和文本解码器组成，不包含独立的文本编码器。

在基础模型方面，\citet{berard2016listen}和\citet{weiss2017sequence}的工作证明了端到端语音翻译的可行性，无需中间转写即可直接从源语言语音生成目标语言文本。这些工作表明端到端S2TT可以达到与级联系统竞争的翻译质量，为后续研究奠定了基础。

在此基础上，研究者们从预训练、知识蒸馏和数据增强等多个角度改进基础编码器-解码器模型，以缓解数据稀缺问题。在预训练方面，\citet{bansal2019pre}利用高资源ASR数据预训练语音编码器，再在低资源S2TT数据上微调，有效提升了低资源条件下的翻译性能。在知识蒸馏方面，\citet{liu2019end}将机器翻译模型作为教师，通过知识蒸馏指导S2TT学生模型的训练，实现了文本翻译知识向语音翻译的迁移。在数据增强方面，\citet{jia2019leveraging}利用语音合成技术将文本翻译数据的源文本合成为语音，构造伪平行数据扩充训练集。GigaST\citep{gigast}构建了万小时规模的伪S2TT语料库以支持大规模训练。BT4ST\citep{fang-feng-2023-back}则利用回译和离散语音单元为S2TT生成伪平行数据，在不依赖转写文本的条件下实现了数据增强。此外，\citet{gaido-etal-2021-ctc}提出在编码器内部通过CTC压缩语音序列长度，降低了语音与文本之间的长度差异。

\subsubsection{并联双编码器}\label{sec:rw_s2tt_parallel}

并联双编码器结构同时包含参数独立的语音编码器和文本编码器，二者并行处理各自模态的输入，输出共同送入共享的文本解码器。训练时，模型同时运行语音翻译和文本翻译两条路径；推理时仅使用语音编码器路径。

T-Modules\citep{duquenne-etal-2022-modules}独立训练语音编码器和文本编码器，通过教师-学生学习将二者的输出对齐到联合嵌入空间，使语音编码器的输出分布逼近文本编码器的输出分布，从而实现了零样本跨模态翻译——即在训练时仅使用文本翻译数据，推理时直接以语音作为输入进行翻译。\citet{ctcmeetsot}提出了孪生并联编码器结构，结合CTC和最优传输进行预训练，通过序列级别的对齐损失从预训练阶段缩小语音和文本之间的模态鸿沟。SeamlessM4T\citep{communication2023seamlessm4t}构建了一站式多语言多模态翻译模型，采用并联的w2v-BERT 2.0\citep{chung2021w2vbert}语音编码器和NLLB\citep{costa2022nllb}文本编码器，通过大规模多任务训练覆盖多达101种语言的语音识别、语音翻译、文本翻译等多种翻译方向，是迄今为止覆盖面最广的统一翻译系统之一。SeamlessM4T v2\citep{communication2023seamlessm4tv2}在此基础上引入了更新的分块注意力机制以支持流式推理。

\subsubsection{串联双编码器}\label{sec:rw_s2tt_serial}

串联双编码器结构将语音编码器的输出送入文本编码器进行进一步处理，再由解码器生成目标文本。这种堆叠结构中，下层语音编码器提取声学特征，上层文本编码器提取语义和翻译特征。根据训练方法的不同，串联双编码器可分为两类：基于机器翻译预训练的迁移学习，以及语音-文本联合预训练。

\textbf{基于机器翻译预训练的迁移学习}\quad 这类方法通常先在机器翻译数据上预训练文本编码器和解码器，再引入语音编码器进行语音翻译微调，将机器翻译知识迁移到语音翻译任务中。早期工作探索了多种串联双编码器的架构设计和跨模态对齐方法。\citet{anastasopoulos2018tied}通过绑定多任务学习，共享编码器并使用双解码器同时训练语音识别和语音翻译任务。Chimera\citep{han-etal-2021-learning}利用wav2vec 2.0作为语音前端，通过共享语义记忆投影模块连接语音表征和文本翻译空间。\citet{dong2021listen}提出声学编码器加语义编码器（以BERT初始化）的架构，通过三重监督实现了声学特征和语义特征的有效解耦。SATE\citep{xu2021sate}构建了堆叠声学-文本编码结构，将语音识别编码器的输出通过适配器连接机器翻译编码器。\citet{liu2020_bridging}提出通过CTC压缩语音编码器输出后送入共享语义编码器，使语义编码器同时处理压缩后的语音特征和文本嵌入，弥合语音和文本之间的模态鸿沟。

后续工作进一步探索了多种训练方法来缩小语音和文本之间的模态鸿沟。\citet{tang-etal-2021-improving}研究了独立的语音编码器和文本编码器之间的底层共享和注意力正则化对齐。XSTNet\citep{ye2021xstnet}通过共享Transformer编码器和解码器，采用跨模态渐进训练促进两种模态的融合。STEMM\citep{fang-etal-2022-STEMM}提出了共享翻译编码器处理语音、文本和混合序列的架构，通过流形混合（Manifold Mixup）在特征空间中缩小模态鸿沟。ConST\citep{ye-etal-2022-cross}通过对比学习拉近句子级别的跨模态表征。M3ST\citep{cheng2023m3st}提出了三层级混合策略，在输入层、编码器层和注意力层分别进行语音-文本混合，充分利用多粒度的跨模态信息。WACO\citep{ouyang-etal-2023-waco}利用词对齐信息进行细粒度的对比学习，在词级别拉近语音和文本的表征。CKD-ST\citep{lei-etal-2023-ckdst}提出了综合知识蒸馏框架，从机器翻译教师模型向语音翻译学生模型进行多层次的知识迁移。CRESS\citep{fang-and-feng-2023-understanding}在语音编码器和文本编码器之上构建共享的翻译编码器和解码器，同时处理语音翻译和机器翻译任务，通过跨模态正则化损失约束两种模态的编码器输出在共享空间中对齐，并采用计划采样（Scheduled Sampling）缓解训练与推理之间的差距。CMOT\citep{zhou-etal-2023-cmot}针对语音和文本特征序列之间的细粒度对齐问题，提出跨模态最优传输方法，利用最优传输理论计算语音帧和文本词元之间的最优匹配方案，实现了特征级别的精确对齐。\citet{zhang2023rethinking}重新审视了语音到文本翻译中的多任务学习策略，通过独立的语音编码器和文本编码器配合长度桥接模块，在多个基准上取得了优异的性能。\citet{han2023tab}系统对比了模态适配和模态正则化两类方法在端到端语音翻译中的表现，为该领域的方法选择提供了实证依据。\citet{zhang2024salign}提出对抗式软对齐方法，通过对抗训练使语音编码器的输出逼近文本编码器的输出分布，实现了更平滑的跨模态对齐。

\textbf{语音-文本联合预训练}\quad 另一类方法不依赖MT预训练，而是同时在语音和文本数据上从头进行联合预训练，使模型在预训练阶段就学习统一的跨模态表示。SLAM\citep{slam}将语音模型和语言模型进行联合预训练，通过在共享的Transformer编码器上同时训练语音和文本任务，在语音翻译等下游任务上取得了显著的效果。mSLAM\citep{mSLAM}将SLAM扩展到多语言场景，覆盖51种语言和43万小时的语音数据，在跨语言语音翻译上表现突出。FUSE\citep{zheng2021fused}提出了融合声学和文本编码的双语预训练方法，在共享编码器中同时处理语音和文本输入，通过多语言多模态预训练增强了语音翻译性能。SpeechUT\citep{zhang2022speechut}通过隐藏单元桥接语音和文本的编解码器，以离散单元作为语音和文本之间的共享中间表示，实现了语音-文本联合表示学习。STPT\citep{tang-etal-2022-unified}设计了部分共享编码器（Partially Shared Encoder）架构，通过模态专属层和共享层的组合进行预训练，在语音翻译和语音识别任务上均取得了优异的性能。


\subsection{语音到语音翻译}\label{sec:rw_encdec_s2st}

基于编码器-解码器的语音到语音翻译（Speech-to-Speech Translation, S2ST）是语音到语音生成领域的核心任务之一，旨在通过端到端模型直接完成从源语言语音到目标语言语音的转换。按照解码策略和模型结构的不同，如图\ref{fig:rw_s2st}所示，现有S2ST方法可分为一阶段S2ST、两阶段S2ST和非自回归语音到语音翻译三类。

% 图：语音到语音翻译流程图
\begin{figure}[htbp]
\centering
\begin{tikzpicture}[node distance=0.3cm and 0.25cm,
    blk/.style={draw, rounded corners=2pt, minimum height=0.7cm,
        font=\small, align=center, thick, fill=#1!12, draw=#1!70!black},
    blk/.default=blue,
]
    % 统一左右输入输出 x 坐标、语音编码器 x、声码器 x
    % 左右留出足够箭头空间，且两侧箭头等长
    \def\xleft{-1.7}
    \def\xright{10.5}
    \def\xenc{0.8}      % 三行语音编码器中心统一
    \def\xvoc{8.2}      % 三行声码器中心统一

    % === (a) 一阶段解码 ===
    \def\ya{0}
    \node[blk=green, minimum width=2.0cm] (a_enc) at (\xenc, \ya) {语音编码器};
    \node[blk=red, minimum width=2.0cm] (a_dec) at (4.4, \ya) {语音解码器};
    \node[blk=teal, minimum width=1.5cm] (a_voc) at (\xvoc, \ya) {声码器};
    \node[font=\small] (a_in) at (\xleft, \ya) {源语音};
    \node[font=\small] (a_out) at (\xright, \ya) {目标语音};
    \draw[arrow2] (a_in) -- (a_enc);
    \draw[arrow2] (a_enc) -- (a_dec);
    \draw[arrow2] (a_dec) -- (a_voc);
    \draw[arrow2] (a_voc) -- (a_out);
    \node[font=\small\bfseries] at (4.4, \ya-0.9) {(a) 一阶段解码};

    % === (b) 两阶段解码 ===
    \def\yb{-2.7}
    \node[blk=green, minimum width=2.0cm] (b_enc) at (\xenc, \yb) {语音编码器};
    \node[blk=orange, minimum width=2.0cm] (b_dec1) at (3.2, \yb) {文本解码器};
    \node[blk=red, minimum width=2.0cm] (b_dec2) at (6.0, \yb) {语音解码器};
    \node[blk=teal, minimum width=1.5cm] (b_voc) at (\xvoc, \yb) {声码器};
    \node[font=\small] (b_in) at (\xleft, \yb) {源语音};
    \node[font=\small] (b_out) at (\xright, \yb) {目标语音};
    \draw[arrow2] (b_in) -- (b_enc);
    \draw[arrow2] (b_enc) -- (b_dec1);
    \draw[arrow2] (b_dec1) -- (b_dec2);
    \draw[arrow2] (b_dec2) -- (b_voc);
    \draw[arrow2] (b_voc) -- (b_out);
    % 标注两个阶段
    \draw[braceup] ([yshift=2pt]b_dec1.north west) -- ([yshift=2pt]b_dec1.north east) node[midway, above=5pt, slbl]{第一阶段};
    \draw[braceup] ([yshift=2pt]b_dec2.north west) -- ([yshift=2pt]b_dec2.north east) node[midway, above=5pt, slbl]{第二阶段};
    % 文本：从文本解码器和语音解码器之间的箭头中点向上引出
    \coordinate (b_mid) at ($(b_dec1.east)!0.5!(b_dec2.west)$);
    \node[above=0.65cm of b_mid, font=\scriptsize] (b_txt) {文本};
    \draw[arrow2] (b_mid) -- (b_txt);
    \node[font=\small\bfseries] at (4.4, \yb-0.9) {(b) 两阶段解码};

    % === (c) 非自回归解码 ===
    \def\yc{-4.8}
    \node[blk=green, minimum width=2.0cm] (c_enc) at (\xenc, \yc) {语音编码器};
    \node[blk=violet, minimum width=2.8cm] (c_dec) at (4.4, \yc) {非自回归解码器};
    \node[blk=teal, minimum width=1.5cm] (c_voc) at (\xvoc, \yc) {声码器};
    \node[font=\small] (c_in) at (\xleft, \yc) {源语音};
    \node[font=\small] (c_out) at (\xright, \yc) {目标语音};
    \draw[arrow2] (c_in) -- (c_enc);
    \draw[arrow2] (c_enc) -- (c_dec);
    \draw[arrow2] (c_dec) -- (c_voc);
    \draw[arrow2] (c_voc) -- (c_out);
    \node[font=\small\bfseries] at (4.4, \yc-0.9) {(c) 非自回归解码};
\end{tikzpicture}
\bicaption{\enspace 语音到语音翻译的三种方法}{\enspace Three approaches for speech-to-speech translation}
\label{fig:rw_s2st}
\end{figure}

\subsubsection{一阶段语音到语音翻译}\label{sec:rw_s2st_1pass}

一阶段S2ST方法通过编码器-解码器模型直接从源语音生成目标语音表征，不依赖显式的文本中间表示。

Translatotron\citep{translatotron}是端到端S2ST的早期代表性工作，采用基于注意力机制的编码器-解码器结构直接从源语音生成目标语音的梅尔谱，证明了无需中间文本即可实现端到端语音翻译的可行性，同时通过辅助的说话人编码器保留了源语音的说话人特征。在离散语音表征方面，\citet{tjandra2019s2st}利用VQ-VAE提取的离散编码进行了跨语言S2ST的初步探索，证明了通过离散中间表示实现语音到语音转换的可行性。UWSpeech\citep{uwspeech}通过跨语言VQ-VAE离散编码建立语音翻译目标，面向无文字语言的语音翻译场景，使得没有书写系统的语言也能进行语音翻译。随后，基于自监督语音模型HuBERT的离散单元成为一阶段S2ST的主流预测目标。S2UT\citep{s2ut}将语音翻译建模为从源语音到目标语音HuBERT离散单元序列的映射，采用Transformer编码器-解码器架构，在Fisher西英数据集上取得了与级联系统可比的翻译质量，奠定了基于离散单元的S2ST研究基础。在此基础上，\citet{lee-etal-2022-textless}进一步实现了完全不依赖文本监督的S2ST系统，仅使用未对齐的语音数据通过挖掘跨语言语音对进行训练。\citet{DBLP:conf/interspeech/PopuriCWPAGHL22}通过自监督预训练编码器和多种数据增强策略（包括语速扰动、噪声添加等）提升了S2UT的翻译质量和鲁棒性。

\subsubsection{两阶段语音到语音翻译}\label{sec:rw_s2st_2pass}

两阶段S2ST方法将翻译过程分解为两个阶段：第一阶段生成目标文本或语义表示，第二阶段基于第一阶段的输出生成目标语音。这种分解策略通过降低单阶段建模的难度来提升翻译质量。

Translatotron 2\citep{translatotron2}是首个采用两阶段架构的端到端S2ST模型，第一阶段由语言解码器生成目标文本，第二阶段由声学合成器基于文本和源语音特征生成目标语音的梅尔谱。这种架构通过将语义翻译和声学合成两个本质不同的子任务解耦，降低了单个模型的建模难度，在翻译质量上达到了与级联系统可比的水平。UnitY\citep{inaguma-etal-2023-unity}采用了类似的两阶段架构，但将第二阶段的生成目标从梅尔谱改为HuBERT离散单元序列，降低了声学建模的复杂度，相比一阶段S2UT取得了2.5至4.2个ASR-BLEU的显著提升。SeamlessM4T\citep{communication2023seamlessm4t}在UnitY的两阶段架构基础上进行了大规模数据扩展，通过百万小时级别的语音数据和多任务训练，构建了覆盖多达101种语言的统一S2ST系统。StreamSpeech\citep{zhang-etal-2024-streamspeech}将两阶段S2ST与同声传译相结合，设计了基于CTC的读写策略来控制翻译延迟，通过多任务学习同时优化语音识别、语音翻译和语音合成目标，实现了同声传译场景下的语音到语音翻译。

\subsubsection{非自回归语音到语音翻译}\label{sec:rw_s2st_nar}

上述一阶段和两阶段S2ST方法均采用自回归解码，对于长语音序列而言解码速度较慢，难以满足实时交互的需求。非自回归（Non-Autoregressive, NAR）方法通过打破逐步解码的依赖，实现并行生成以大幅提升推理速度。文本领域的非自回归翻译研究为NAR S2ST提供了重要的技术基础，代表性方法包括NAT\citep{gu2018nat}、CMLM\citep{cmlm}、GLAT\citep{qian-etal-2021-glancing}、DA-Transformer\citep{huang2022DATransformer}和CTC\citep{ctc}等。

在非自回归语音到语音翻译方面，TranSpeech\citep{huang2023transpeech}是该方向的早期代表性工作，采用CMLM进行并行解码。TranSpeech首先通过双边扰动（Bilateral Perturbation）对语音进行风格归一化和信息增强，降低离散单元的声学多峰性，使相同语义内容对应更确定的离散表示，然后利用CMLM在多轮迭代中逐步生成完整序列，在保持翻译质量的同时实现了21.4倍的解码加速。DiffS2UT\citep{diffs2ut}将扩散模型引入S2ST任务，在离散单元的连续嵌入空间上执行扩散过程，通过迭代去噪实现非自回归生成。\citet{wu-2023-duplex}提出了双工扩散模型用于S2ST，采用双向扩散过程分别建模源到目标和目标到源的翻译，进一步提升了扩散方法在该任务上的表现。DiffNorm\citep{tan2024diffnorm}从简化训练目标的角度出发，通过扩散归一化将离散单元的不规则分布转换为近似高斯分布，使非自回归模型更容易拟合目标分布，降低了训练难度。NAST-S2x\citep{ma2024nasts2x}则将非自回归方法与同声传译相结合，基于CTC实现了端到端的非自回归同时语音到语音翻译，在翻译延迟和翻译质量上取得了良好的平衡。

尽管非自回归方法在解码速度上取得了显著提升，但其翻译质量相比先进的自回归模型仍有差距。如何在解码效率和翻译质量之间取得更好的平衡，是当前S2ST研究面临的核心挑战。同时，现有的S2ST系统大多需要从头训练，难以复用已有的预训练S2TT和TTS模型。


%===========================================================================
% 2.4 基于仅解码器模型的语音到语音生成
%===========================================================================
\section{基于大语言模型的语音对话}\label{sec:rw_deconly_all}

随着大语言模型的兴起，语音到语音生成的研究范围从翻译拓展到更加通用的语音对话，仅解码器架构成为这一方向的主流范式。本节首先介绍仅解码器架构的基础知识，随后分别介绍语音理解大模型和语音对话大模型的相关研究。

\subsection{仅解码器架构基础}\label{sec:rw_deconly}

随着大语言模型的兴起，仅解码器（Decoder-only）架构逐渐取代编码器-解码器架构，成为序列生成任务的主流范式。与编码器-解码器架构不同，如图\ref{fig:rw_decoderonly}所示，仅解码器架构移除了编码器和交叉注意力模块，仅保留了因果自注意力和前馈神经网络的堆叠。

% 图：仅解码器架构示意图
\begin{figure}[htbp]
\centering
\begin{tikzpicture}[node distance=0.35cm,
    block/.style={draw, rounded corners=2pt, minimum width=2.8cm, minimum height=0.6cm,
        font=\small, align=center, thick, fill=#1!12, draw=#1!70!black},
    block/.default=blue]

    \node[block=green, minimum width=3.0cm] (input) at (0, 0) {输入嵌入 + 位置编码};
    \node[block=orange, above=of input, minimum width=3.0cm] (sa) {因果自注意力};
    \node[block=orange, above=of sa, minimum width=3.0cm] (ffn) {前馈网络};
    \node[block=gray, above=0.5cm of ffn, minimum width=3.0cm] (linear) {线性层 + Softmax};

    \draw[arrow2] (input) -- (sa);
    \draw[arrow2] (sa) -- (ffn);
    \draw[arrow2] (ffn) -- (linear);

    \begin{scope}[on background layer]
        \node[dbox=orange!40, fit=(sa)(ffn), inner sep=6pt] (dec_layer) {};
    \end{scope}
    \node[right=4pt of dec_layer.east, font=\scriptsize, orange!70] {$\times N$};

    \node[below=0.25cm of input, font=\small] {输入序列};
    \node[above=0.25cm of linear, font=\small] {下一词元预测};
\end{tikzpicture}
\bicaption{\enspace 仅解码器架构}{\enspace Decoder-only architecture}
\label{fig:rw_decoderonly}
\end{figure}

GPT\citep{gpt1}是基于仅解码器Transformer架构的生成式预训练语言模型的奠基性工作。GPT以标准语言建模为训练目标，即自回归地预测下一个词元，在大规模无标注文本上进行预训练，然后通过在下游任务上微调来适配具体应用。GPT提出的"预训练-微调"范式证明了单向语言模型可以学习到丰富的语言表示，为后续大语言模型的发展奠定了基础。GPT-2\citep{radford2019language}将模型规模扩展到15亿参数，并在WebText等大规模语料上训练，展示了大规模语言模型在不进行任何微调的情况下完成多种任务的零样本能力。GPT-3\citep{brown2020gpt3}进一步将参数量扩大至1750亿，系统性地验证了上下文学习（In-Context Learning）能力——仅通过在提示中提供少量示例，模型即可完成翻译、问答和推理等多种任务，无需更新模型参数。这些工作共同确立了仅解码器架构在语言生成领域的主导地位。

后续的大语言模型延续了仅解码器架构和自回归语言建模的技术路线，并在此基础上不断扩大模型规模和训练数据。LLaMA\citep{touvron2023llama}系列是推动开源大语言模型发展的重要力量，LLaMA初始版本覆盖7B至65B参数规模，LLaMA 2\citep{touvron2023llama2}进一步扩展至70B并引入基于人类反馈的强化学习（RLHF），LLaMA 3\citep{dubey2024llama3}则达到8B/70B/405B参数，训练数据规模超过15万亿词元。Qwen2\citep{yang2024qwen2}和Qwen2.5\citep{yang2025qwen25}是阿里通义千问系列的代表，在中英双语和多语言任务上表现突出。

仅解码器架构的训练目标为标准语言建模，即最小化下一个词元的负对数似然：
\begin{equation}
\mathcal{L}(\theta)=-\sum_{i=1}^{M}\log P_\theta(y_i|Y_{<i}),
\end{equation}
其中$Y_{<i}$为位置$i$之前的序列前缀。与编码器-解码器架构的训练目标相比，仅解码器架构的条件中不包含独立的源序列$X$，而是将所有输入（包括指令、上下文等）统一为序列前缀，以自回归方式生成后续内容。推理时，模型通过提示（Prompt）引导生成，将输入信息编码为提示序列后自回归地生成目标序列。

\textbf{语音语言建模}\quad 在大语言模型的影响下，研究者开始探索将语言建模的思想应用于语音模态，即在离散语音词元上训练自回归语言模型，使模型能够以无文本的方式直接建模和生成语音。GSLM\citep{lakhotia2021generative}是语音语言建模（Spoken Language Modeling）的早期代表性工作，在HuBERT离散单元上训练自回归Transformer语言模型，实现了无需文本监督的语音生成和语音续写。AudioLM\citep{borsos2023audiolm}提出了层级化的自回归生成框架：首先利用语义词元（来自W2V-BERT）建模语音的语义内容，再通过声学词元（来自SoundStream）恢复声学细节和音质，这种从粗到细的生成策略在语音续写和音频生成上取得了高度自然的效果。VALL-E\citep{wang-etal-2023-valle}将零样本语音合成建模为语言模型任务，以3秒语音提示即可生成保持说话人音色的合成语音。VioLA\citep{wang2023viola}统一编解码框架处理语音识别、语音翻译和语音合成等多种任务，验证了单一语言模型架构在多种语音任务上的适用性。这些工作为后续基于大语言模型的语音到语音生成奠定了基础。


\subsection{语音理解大模型}\label{sec:rw_deconly_s2t}

基于仅解码器模型的语音到文本生成旨在利用大语言模型的文本理解和生成能力处理语音输入。如图\ref{fig:rw_s2t_llm}所示，其典型架构由语音编码器、适配器（Adaptor）和大语言模型三部分组成：语音编码器（如Whisper、Wav2Vec 2.0等）提取连续语音特征，适配器将语音特征映射到大语言模型的文本嵌入空间中，大语言模型基于映射后的语音表征生成文本输出。由于语音编码器和大语言模型均可灵活替换（如语音编码器可选用Whisper或Wav2Vec 2.0，大语言模型可选用LLaMA或Qwen等），适配器的设计成为该类工作的核心差异所在。按照适配器类型的不同，现有工作可分为以下几类。

% 图：语音到文本生成架构图
\begin{figure}[htbp]
\centering
\begin{tikzpicture}[node distance=0.3cm and 0.4cm,
    blk/.style={draw, rounded corners=2pt, minimum height=0.8cm,
        font=\small, align=center, thick, fill=#1!12, draw=#1!70!black},
    blk/.default=blue,
]
    \node[blk=green, minimum width=2.2cm] (senc) {语音编码器\\{\scriptsize (Whisper 等)}};
    \node[blk=orange, right=1.3cm of senc, minimum width=2.0cm] (adapt) {适配器\\{\scriptsize (MLP / Q-Former / CTC)}};
    \node[blk=red, right=0.8cm of adapt, minimum width=2.5cm] (llm) {大语言模型\\{\scriptsize (LLaMA / Qwen 等)}};

    \node[left=0.4cm of senc, font=\small] (in) {语音};
    \node[right=0.4cm of llm, font=\small] (out) {文本};

    \draw[arrow] (in) -- (senc);
    \draw[arrow] (senc) -- node[above=1pt, slbl, fill=white, inner sep=1pt]{语音表征} (adapt);
    \draw[arrow] (adapt) -- node[above=1pt, slbl, fill=white, inner sep=1pt]{嵌入} (llm);
    \draw[arrow] (llm) -- (out);
\end{tikzpicture}
\bicaption{\enspace 基于大语言模型的语音到文本生成架构}{\enspace Architecture of LLM-based speech-to-text generation}
\label{fig:rw_s2t_llm}
\end{figure}

\textbf{多层感知机适配器}\quad 多层感知机（Multi-Layer Perceptron, MLP）是最简洁的适配器设计，通过一个或多个线性层将语音编码器的输出投影到大语言模型的嵌入空间。Qwen-Audio\citep{Qwen-Audio}将Whisper编码器通过投影层连接Qwen-7B大语言模型，构建了支持多种音频理解任务的统一模型。Qwen2-Audio\citep{Qwen2-Audio}在此基础上采用Whisper-large-v3编码器，并通过直接偏好优化（Direct Preference Optimization, DPO）进一步提升了指令跟随能力。SLAM-ASR\citep{slam-asr}构建了模块化框架，以线性投影器连接语音编码器和大语言模型，在语音识别任务上取得了优异的结果。\citet{fathullah2024prompting}通过帧堆叠和线性投影将CTC预训练Conformer编码器的输出映射到大语言模型输入空间。LLAST\citep{chen-etal-2024-llast}采用三层MLP适配器连接Whisper编码器和大语言模型，通过ASR增强训练实现了端到端语音翻译。

\textbf{Q-Former适配器}\quad Q-Former是一种基于可学习查询和交叉注意力的适配器结构，最初由BLIP-2\citep{li2023blip2}提出并应用于视觉-语言对齐。其核心思想是使用一组固定数量的可学习查询向量，通过交叉注意力从编码器输出中提取信息，从而将变长的编码器输出压缩为固定长度的表示。SALMONN\citep{tang2024salmonn}采用双编码器架构（Whisper和BEATs），使用窗口级Q-Former将两个编码器的输出融合后连接至Vicuna-13B大语言模型，实现了听觉感知和理解的统一。\citet{yu2024connecting}系统地对比了全连接层、交叉注意力和Q-Former三种连接方式，发现Q-Former在固定长度压缩上表现优异但对变长输入适应性不足，进而提出了段级Q-Former，将语音编码器输出分段后分别进行Q-Former压缩再拼接，以更好地适配语音的可变长度特性。

\textbf{CTC适配器}\quad CTC适配器利用CTC的对齐信息对语音编码器的输出进行压缩和下采样，将变长的语音特征序列转换为与文本长度更接近的紧凑表示后送入大语言模型。与MLP和Q-Former相比，CTC适配器能够利用语音内容信息自适应地确定压缩位置，更有效地去除语音中的冗余帧（如静音和重复帧），使得压缩后的表示序列长度接近文本词元长度。Speech-LLaMA\citep{wu2023decoder}采用CTC压缩将变长语音编码器输出压缩为定长序列，在低资源条件下的语音识别任务上表现优异。\citet{hono-etal-2024-integrating}系统对比了卷积下采样、CTC移除（丢弃CTC空白帧）和CTC平均（将相邻同标签帧取平均）三种桥接方式，发现CTC平均在语音识别任务上表现最优，因为它在有效压缩序列的同时保留了更多的声学信息。

除上述适配器外，部分工作采用了混合策略。WavLLM\citep{hu2024wavllm}结合双编码器（Whisper和WavLM）与卷积-瓶颈-线性复合适配器，并引入提示感知的LoRA实现了对不同语音任务的自适应处理。


\subsection{语音对话大模型}\label{sec:rw_deconly_s2s}

基于仅解码器模型的语音到语音生成是该领域的前沿方向，其核心目标是使大语言模型能够直接以语音为输入和输出进行交互，绕开传统的ASR加LLM加TTS级联流程。实现这一目标面临多维度的挑战：首先，赋予大语言模型语音生成能力需要大规模的语音训练数据，而大规模语音训练可能导致模型原有的文本理解和推理能力出现灾难性遗忘；其次，语音对话场景对响应延迟有严格要求，模型需要在生成文本的同时支持语音的流式输出，而非等待文本完全生成后再合成语音；最后，语音解码器的设计需要在生成质量和流式输出之间取得平衡。围绕这些挑战，研究者沿两条技术路线展开探索：一条是将语音词元直接融入大语言模型的词表进行联合训练，使模型内在地具备语音生成能力；另一条是保持大语言模型的文本能力不变，在其外部挂接独立的语音解码器来完成语音生成。如图\ref{fig:rw_s2s_llm}所示，这两条路线分别对应原生式和组合式两类方法。

% 图：语音到语音生成架构图（原生式 vs 组合式）
\begin{figure}[htbp]
\centering
\begin{tikzpicture}[node distance=0.3cm and 0.35cm,
    blk/.style={draw, rounded corners=2pt, minimum height=0.7cm,
        font=\small, align=center, thick, fill=#1!12, draw=#1!70!black},
    blk/.default=blue,
]
    % 固定左右输入输出 x 坐标，两子图共用
    \def\xleft{-1.0}
    \def\xright{9.8}

    % === (a) 原生式 ===
    \def\ya{0}
    \node[blk=violet, minimum width=5.0cm, minimum height=1.2cm] (native) at (4.4, \ya) {
        语音-语言模型\\
        {\scriptsize 统一词表：文本词元 + 语音词元}
    };
    \node[font=\small] at (\xleft, \ya) {语音};
    \node[font=\small] at (\xright, \ya) {语音};
    \draw[arrow] (\xleft+0.4, \ya) -- (native);
    \draw[arrow] (native) -- (\xright-0.4, \ya);
    \node[font=\small\bfseries] at (4.4, \ya-1.2) {(a) 原生式};

    % === (b) 组合式 ===
    \def\yb{-3.2}
    \node[blk=green, minimum width=1.8cm] (b_senc) at (1.0, \yb) {语音编码器};
    \node[blk=orange, minimum width=1.5cm] (b_adapt) at (3.3, \yb) {适配器};
    \node[blk=red, minimum width=2.0cm] (b_llm) at (5.5, \yb) {大语言模型};
    \node[blk=blue, minimum width=1.8cm] (b_sdec) at (7.9, \yb) {语音解码器};

    \node[font=\small] at (\xleft, \yb) {语音};
    \node[font=\small] at (\xright, \yb) {语音};

    \draw[arrow2] (\xleft+0.4, \yb) -- (b_senc);
    \draw[arrow2] (b_senc) -- (b_adapt);
    \draw[arrow2] (b_adapt) -- (b_llm);
    \draw[arrow2] (b_llm) -- (b_sdec);
    \draw[arrow2] (b_sdec) -- (\xright-0.4, \yb);

    % 文本输出
    \node[above=0.4cm of b_llm, font=\scriptsize] (txt_out) {文本};
    \draw[arrow2] (b_llm.north) -- (txt_out);

    \node[font=\small\bfseries] at (4.4, \yb-0.9) {(b) 组合式};
\end{tikzpicture}
\bicaption{\enspace 语音到语音生成的两种架构：原生式与组合式}{\enspace Two architectures for speech-to-speech generation: native vs.\ modular}
\label{fig:rw_s2s_llm}
\end{figure}

\subsubsection{原生式语音大模型}\label{sec:rw_s2s_native}

原生式语音大模型在大语言模型内部统一建模文本和语音词元，使模型具备直接生成语音的能力。

最早将文本大语言模型扩展到语音模态的工作采用离散词元同时表示语音的输入和输出。TWIST\citep{twist}利用文本大语言模型（如LLaMA）的参数初始化语音语言模型，将文本词表替换为语音离散词元词表后继续在语音数据上训练，验证了文本预训练知识向语音模态迁移的可行性，发现文本预训练能够显著提升语音语言模型的语义理解能力。SpeechGPT\citep{zhang-etal-2023-speechgpt}将HuBERT离散语音词元扩展到LLaMA的词表中，通过跨模态指令微调使模型具备了语音输入和语音输出的对话能力，是跨模态语音对话大语言模型的早期代表性工作。SpeechGPT采用三阶段训练策略：先进行模态适配预训练，再进行跨模态指令微调，最后通过思维链微调提升推理能力。AudioPaLM\citep{rubenstein2023audiopalm}将PaLM-2大语言模型与AudioLM的语音生成能力相融合，保留PaLM-2的语言知识同时赋予其处理音频的能力，支持语音识别、语音翻译和语音到语音翻译等多种任务。

离散语音词元与文本词元之间存在粒度差异：相同语义内容的语音词元序列通常远长于对应的文本词元序列。为解决这一问题，研究者探索了多种语音-文本对齐方案。SpiRit-LM\citep{nguyen2024spirit}通过在训练数据中交错（Interleave）排列语音词元和文本词元，使模型学习两种模态之间的对齐关系。GLM-4-Voice\citep{glm4voice}同样采用语音-文本交错的训练策略，并通过监督微调增强模型的对话能力。SpeechGPT-Gen\citep{zhang2024speechgptgen}引入链式信息生成策略，通过分组对齐降低语音和文本之间的粒度差异。IntrinsicVoice\citep{zhang2024intrinsicvoice}设计了GroupFormer模块对语音词元进行分组压缩，同样实现了语音-文本粒度的匹配。

上述离散输入方案在将语音转换为离散词元的过程中不可避免地损失部分信息。为此，一些工作采用连续语音特征作为输入以保留更完整的声学信息，同时使用离散词元作为输出以利用自回归语言模型的生成能力。Mini-Omni\citep{xie2024mini}是较早的开源实时语音对话模型，采用Whisper编码器提取连续语音特征作为输入，通过并行的文本头和语音头使大语言模型能够同时生成文本词元和语音词元，实现了端到端的语音到语音对话。Kimi-Audio\citep{moonshot2025kimiaudio}构建了通用的音频基础模型，采用大规模音频-文本交错预训练策略，以连续语音特征为输入、离散音频词元为输出，同时支持语音理解、音频理解、语音对话和音频生成等多种任务。

传统的语音对话遵循轮替（Turn-Taking）模式，即一方说完另一方再回复，这种模式无法处理打断、插话和同时说话等自然对话中的常见现象。全双工（Full-Duplex）对话使模型能够在说话的同时监听用户输入，实现更自然的交互体验。dGSLM\citep{nguyen2022dgslm}是全双工方向的早期探索性工作，提出了双通道对话语言模型，使用两个独立的Transformer分别建模对话双方的语音流，并通过交叉注意力实现通道间的交互。虽非仅解码器架构，但dGSLM验证了建模双向语音流的可行性，奠定了全双工语音对话的研究方向。Moshi\citep{kyutai2024moshi}是实时全双工语音对话大语言模型的代表性工作，其核心是一个自回归的Helium语言模型，通过Inner Monologue机制同时生成文本词元和Mimi编解码器的多层级语音词元，利用深度RQ-Transformer架构并行预测多个RVQ层级，在仅200毫秒理论延迟的条件下实现了流畅的全双工对话。工业界也在积极推进语音到语音生成技术，GPT-4o\citep{gpt4o}是OpenAI发布的原生多模态大模型，能够直接以语音为输入和输出进行交互，在语音响应的自然度和情感表现力上达到了接近人类对话的水平。Gemini Live\citep{google2024gemini}是Google推出的实时语音对话系统，支持自然的多轮语音对话和打断功能。

\subsubsection{组合式语音大模型}\label{sec:rw_s2s_composite}

组合式语音大模型将语音理解和语音生成功能解耦：大语言模型负责文本理解和生成，外接的语音解码器将大语言模型的文本输出或隐藏状态转换为语音。相比原生式方法，组合式架构避免了在大语言模型内部引入大量语音数据所带来的文本能力遗忘问题，且模块化的设计降低了训练成本。

Freeze-Omni\citep{xiong2024freeze}提出在语音对话训练过程中完全冻结大语言模型的参数，仅训练轻量的语音编码器和语音解码器，最大程度地保留了大语言模型原有的文本理解和推理能力，有效避免了灾难性遗忘问题。Minmo\citep{chen2025minmo}基于约8B参数的大语言模型搭配CosyVoice 2语音解码器，通过多阶段训练策略对齐语音编码器、大语言模型和语音解码器，构建了支持无缝语音对话的完整系统。Qwen2.5-Omni\citep{alibaba2025qwen25omni}提出了Thinker-Talker架构，其中Thinker组件负责接收多模态输入并进行文本推理，Talker组件基于Thinker的隐藏状态以流式方式生成语音词元，将文本推理和语音生成功能清晰解耦，在保证文本生成质量的同时实现了高质量的语音输出。VocalNet\citep{vocalnet2025}引入多词元预测策略，在每个解码步骤同时预测多个语音词元以加速语音生成，在解码效率上取得了显著提升。

% \subsubsection{全模态大模型}\label{sec:rw_s2s_omni}
%
% 部分研究将语音模态纳入更广泛的全模态（Omni-modal）框架中，构建能够同时处理语音、视觉和文本等多种模态的统一模型。AnyGPT\citep{zhan2024anygpt}采用统一的离散token表示文本、语音、图像和音乐等多种模态，在单一的自回归Transformer中实现了任意模态之间的转换。VITA\citep{fu2024vita}构建了开源的交互式全模态大语言模型，同时支持视频、图像、文本和音频的输入以及文本和语音的输出，通过状态token实现了非唤醒条件下的交互。EMOVA\citep{chen2025emova}设计了语义-声学解耦分词器，将语音的语义内容和情感特征分别编码，使模型在全模态交互中能够生成包含丰富情感（如高兴、悲伤、愤怒等）的语音响应。这些工作代表了从单模态向全模态统一建模演进的趋势。

综合来看，原生式方法在语音对话的自然度上具有优势，但需要海量语音数据进行训练，且容易遗忘大语言模型原有的文本能力。组合式方法通过模块化设计在训练效率和文本能力保持上具有优势，但语音解码器的质量和流式生成能力仍有待提升。


%===========================================================================
% 2.5 研究现状综合分析
%===========================================================================
\section{研究现状综合分析}\label{sec:rw_analysis}

本章系统梳理了语音到语音生成领域的研究现状，涵盖了从编码器-解码器架构到仅解码器架构的两条技术发展脉络。在编码器-解码器范式下，研究从语音到文本翻译（S2TT）的模型结构设计、跨模态对齐和训练策略出发，逐步拓展到语音到语音翻译（S2ST）的一阶段、两阶段和非自回归方法。在仅解码器范式下，研究从利用适配器赋予大语言模型语音理解能力的语音到文本生成出发，进一步发展为同时具备语音理解和生成能力的语音到语音大语言模型，包括原生式和组合式两类技术路线。这两条技术发展脉络交汇于共同的目标：实现高质量、低延迟的端到端语音到语音生成。

在对现有研究的全面分析中，可以发现以下四个尚未解决的关键科学问题，这些问题引出了本文的五项研究工作。

现有的端到端S2ST模型大多需要从头训练，难以复用已有的预训练S2TT模型和TTS模型。其根本原因在于不同模块之间的接口不兼容，包括模型结构、特征空间和输入输出词表的差异。如何实现预训练模块的灵活组合，以降低训练成本并提升系统性能，是S2ST系统构建中的重要问题。针对\textit{预训练模型的组合与利用}这一问题，本文第\ref{chap:comspeech}章提出基于模块化组合的端到端语音到语音翻译模型ComSpeech，通过CTC词表适配器实现S2TT模型和TTS模型的自由组合，在低训练成本下构建了高质量的S2ST系统。

在编码器-解码器范式的S2ST研究中，自回归模型翻译质量高但解码速度慢，而非自回归模型虽然大幅提升了解码速度，但翻译质量相比自回归模型仍有较大差距。如何在翻译质量和解码效率之间取得更好的平衡，是S2ST研究面临的核心挑战。针对这一\textit{建模能力与解码效率的矛盾}，本文第\ref{chap:daspeech}章提出基于有向无环图的语音到语音翻译模型DASpeech，基于有向无环图（DAG）构建高质量非自回归语音到语音翻译模型，在保持接近自回归模型翻译质量的同时实现了大幅度的解码加速。

在仅解码器范式下，如何以低成本赋予大语言模型语音理解和语音生成能力，同时实现文本与语音的同步流式生成以保证低延迟，是语音到语音大语言模型研究面临的重要挑战。原生式方法需要海量语音数据从头训练，成本高昂；而组合式方法中，非自回归语音解码器虽然速度快但生成语音的自然度不足，自回归语音解码器质量更高但因串行解码的特性难以实现流式输出。针对\textit{多模态序列的同步流式生成}这一问题，本文第\ref{chap:llama-omni}章提出基于非自回归流式解码的实时语音对话大模型LLaMA-Omni，以低训练成本构建组合式语音对话大模型，通过非自回归语音解码器实现低延迟的语音响应；第\ref{chap:llama-omni2}章进一步提出基于自回归流式解码的高质量语音对话大模型LLaMA-Omni 2，采用自回归流式语音解码器，在大语言模型生成文本的同时实时合成语音，在准确性和自然度上取得了显著提升。

上述语音对话模型均面向基于轮次的对话场景，依赖外部的语音活动检测模块来判断用户发言边界，无法支持打断、停顿等灵活的交互模式。在真实的人际对话中，双方的发言往往存在重叠和交替，对话时机的把控依赖于对语义内容的深层理解，这远超声学层面的语音活动检测所能提供的能力。如何使模型在持续接收音频输入的同时自主决定回复时机，实现全双工语音对话，是迈向自然人机对话的关键挑战。针对\textit{双向语音流的统一序列建模}这一问题，本文第\ref{chap:full-duplex}章提出基于多信道交织的全双工语音对话模型LLaMA-Omni-Duplex，将用户语音、回复文本和回复语音三个信道按固定比例统一建模，通过文本信道中的状态标记实现端到端的时机决策，无需引入额外的分类模块。

\section{本章小结}

本章全面调研了语音到语音生成领域的研究现状。首先介绍了三种建模范式和语音表征的相关研究。随后，从编码器-解码器和仅解码器两种架构出发，分别梳理了语音到文本生成和语音到语音生成的技术发展脉络。最后，通过综合分析指出了当前研究中存在的四个关键科学问题，明确了本文五项研究工作的出发点。下一章将介绍本文的第一项工作——ComSpeech，一种基于模块化组合的端到端语音到语音翻译模型。

